% !TeX program = xelatex
% SPDX-License-Identifier: MIT
% Copyright (c) 2026 Ziyu Peng and 刘毅涵
% Author: https://github.com/pzy54154631-hub
% Author: https://pzy54154631-hub.github.io/liuyihan/
% 教学模型；所有参数均为自设，不代表真实市场。
\documentclass[UTF8,a4paper,fontset=fandol]{ctexart}
\usepackage[margin=2.5cm]{geometry}
\usepackage{amsmath,amssymb,pgfplots}
\pgfplotsset{compat=1.18}
\usepackage[hidelinks]{hyperref}
\title{经济学模型：预算、选择与市场}
\author{\href{https://github.com/pzy54154631-hub}{\textit{Ziyu Peng}}, University of Colorado Boulder,
  \\[0.1em]College of Arts and Sciences, Statistics and Data Science
  \\[0.1em]Boulder, CO 80309, USA
  \\[0.1em]\href{mailto:Ziyu.PengSr@colorado.edu}{Ziyu.PengSr@colorado.edu}
  \\[0.35em]
  \href{https://pzy54154631-hub.github.io/liuyihan/}{刘毅涵}，暨南大学经济学院}
\date{}
\begin{document}
\maketitle
\section{消费者选择}
设两种商品数量为 $x,y>0$，价格 $p_x,p_y>0$，
收入 $m>0$，偏好参数 $0<\alpha<1$。采用对数效用：
\begin{equation}
  \max_{x,y>0}\; \alpha\ln x+(1-\alpha)\ln y
  \quad\text{s.t.}\quad p_xx+p_yy\le m.
  \label{eq:choice}
\end{equation}
因效用对两种商品均严格递增，最优选择用尽预算。
将约束写成等式后，拉格朗日函数及一阶条件为
\begin{align}
 \mathcal L&=\alpha\ln x+(1-\alpha)\ln y
           +\lambda(m-p_xx-p_yy),\\
 \frac{\partial\mathcal L}{\partial x}
   &=\frac{\alpha}{x}-\lambda p_x=0,\\
 \frac{\partial\mathcal L}{\partial y}
   &=\frac{1-\alpha}{y}-\lambda p_y=0,\\
 p_xx+p_yy&=m.
\end{align}
得到唯一最优解
\[
 x^*=\frac{\alpha m}{p_x},\qquad
 y^*=\frac{(1-\alpha)m}{p_y}.
\]
例如 $m=100,p_x=5,p_y=10,\alpha=0.4$，
有 $x^*=8,y^*=6$，预算核对为 $5\times8+10\times6=100$。
这些值仅是教学设定。

\clearpage
\section{供需示意图}
下面是与上例独立的线性市场模型：
\[
 p_D(q)=12-0.5q,\qquad p_S(q)=2+0.5q.
\]
联立后得到 $q^*=10,p^*=7$。图中横轴是数量，纵轴是价格。
\begin{figure}[htbp]
\centering
\begin{tikzpicture}
\begin{axis}[
 width=0.8\linewidth,height=7cm,
 axis lines=left,xmin=0,xmax=20,ymin=0,ymax=13,
 xlabel={数量 $q$},ylabel={价格 $p$},
 xtick={0,5,10,15,20},ytick={0,2,7,12},
 legend style={at={(0.98,0.98)},anchor=north east},
 samples=2]
\addplot[blue,thick,domain=0:20]{12-0.5*x};
\addlegendentry{需求}
\addplot[red,dashed,thick,domain=0:20]{2+0.5*x};
\addlegendentry{供给}
\addplot[gray,densely dotted,forget plot]
 coordinates {(0,7) (10,7) (10,0)};
\addplot[black,only marks,mark=*,forget plot]
 coordinates {(10,7)};
\node[above left] at (axis cs:10,7) {$E(10,7)$};
\end{axis}
\end{tikzpicture}
\caption{自设线性供需模型；曲线不来自市场数据。}
\label{fig:supply-demand}
\end{figure}
如图~\ref{fig:supply-demand}，虚线投影帮助读出均衡坐标。
\end{document}
